\cotefile{136}
\scanpage{282}
\mspage{- 35 -}\hfill IX
% archive sheet number: 282
% mounting strip and printed label: 353

On a un isomorphisme canonique (compatible aux opérations de $\handwritten{\pi}$) :
\[
H^{1}(A_{\overline{K}},\mathbb{Z}_{\handwritten{\ell}})\simeq\operatorname{Hom}(T_{\handwritten{\ell}}(A(\overline{K})),\mathbb{Z}_{\handwritten{\ell}})\ ,\tag*{(*)}
\]
qui montre que l'assertion 3.7 équivaut à la même assertion pour les opérations de $I$ sur le premier membre de (*). Or cette assertion a été prouvée dans III \qquad (en utilisant la résolution des singularités des schémas excellents de dimension 2).

\uline{Corollaire} 3.8. \uline{Les conditions équivalentes} (i) \uline{à} (iv) \uline{de} 3.5 \uline{équivalent aussi à la condition suivante} :

(vi) \uline{L'action de} $I$ \uline{sur} $U=T_{\handwritten{\ell}}(A(\overline{K}))$ \uline{est unipotente}.

En effet, pour prouver que (v) implique (vi), on note que si un groupe $I$ opérant de façon unipotente sur un vectoriel $E$ de dimension finie sur un corps de caractéristique nulle (en l'occurrence $\mathbb{Q}_{\handwritten{\ell}}$) admet un sous-groupe d'indice fini $I'$ opérant avec l'échelon $m$ (en l'occurrence $m=2$), alors $I$ lui-même opère avec l'échelon $m$. En effet, $I'$ étant d'indice fini dans $I$, le groupe algébrique $G'\handwritten{\subset}\text{\uline{Gl}}(E)$ qu'il définit est d'indice fini dans celui, $G$, défini par $I$, or $G$ étant unipotent est connexe (caractéristique nulle !), donc $G=G'$, donc $I$ et $I'$ ont même ``échelon d'unipotence''.

On obtient également une réciproque à 3.3 :

\uline{Corollaire} 3.9. \uline{Les notations sont celles de} 3.2. \uline{Les conditions} (i) \uline{et} (ii) \uline{de} 3.2 \uline{équivalent aussi à la condition suivante}\handwritten{$\unsure{\backslash}$} :

(iii) \uline{Pour toute extension étale} $K'$ \uline{de} $K$, \uline{désignant par} $S'$ \uline{le normalisé de} $S$ \uline{dans} $K'$, \uline{par} $A'$ \uline{le modèle de Néron de} $A_{K'}=A\otimes_{S}K'=A_{K}\otimes_{K}K'$, \uline{le morphisme canonique}
\scanpage{283}
% archive sheet number: 283
\diagram{\begin{tikzcd}[ampersand replacement=\&,column sep=small]
\cdots P \arrow[r] \& C\otimes C' \arrow[r,shift left=.6ex] \arrow[r,shift right=.6ex] \& C\otimes C'\otimes\unsure{s|S}
\end{tikzcd}}

\diagram{\begin{tikzcd}[ampersand replacement=\&,column sep=small]
V(P) \& V(C\otimes\deleted{\unsure{s\otimes|C}}C')\arrow[l]\arrow[dl] \& V(C)\times V(C')\times V(s)\arrow[l,shift left=.6ex]\arrow[l,shift right=.6ex]\\
V(C'')\&\&
\end{tikzcd}}

\[
\begin{aligned}
P&=\operatorname{Ker}\bigl(C\otimes C'\rightrightarrows C\otimes\unsure{s\otimes|\otimes}C'\bigr)\deleted{\otimes s\,)}\\
P\deleted{\unsure{\otimes x|\otimes s|\otimes}}&\longrightarrow P\otimes s=\operatorname{Ker}\bigl(\underline{C\otimes C'}\otimes s\unsure{=|\rightrightarrows}\underline{C\otimes s\otimes C'}\otimes s\bigr)
\end{aligned}
\]
\[
\begin{aligned}
P&\longrightarrow C\otimes C'\\
P\otimes s&\longrightarrow C\otimes C'\otimes s
\end{aligned}
\]

\diagram{\begin{tikzcd}[ampersand replacement=\&,column sep=large]
P\arrow[d,"?"']\arrow[r,phantom,"{\overset{i}{\subset}}" description] \& C\otimes C'\arrow[d,shift right=.8ex,"\alpha"']\arrow[d,shift left=.8ex,"\beta"]\\
P\otimes s\arrow[r,phantom,"{\overset{i\otimes\mathrm{id}_{s}}{\subset}}" description] \& C\otimes C'\otimes s
\end{tikzcd}}
\[
\alpha i=\beta i\overset{?}{=}(i\otimes\mathrm{id}_{s})\circ\left(\underset{?}{\Delta_{P}}\right)
\]
\[
\begin{aligned}
x&\in C\otimes C' &\alpha(x)&=\beta(x)\\
\deleted{x}\,\unsure{y|a|s}&\in s
\end{aligned}
\]
\[
\alpha(x)=\beta(x)\ \text{\unsure{est|$\longrightarrow$} dans }\deleted{\unsure{\alpha}}\operatorname{Im}(x\otimes s\longrightarrow C\otimes C'\otimes s)
\qquad
\]
\[
x\otimes s
\]

\diagram{\begin{tikzcd}[ampersand replacement=\&,column sep=tiny,row sep=small]
\&\&C''\arrow[dl]\arrow[d]\arrow[dr]\&\\
\cdots P_{3}\arrow[r,phantom,"\subset"]\arrow[d] \& P_{2}\arrow[r,phantom,"\subset"]\arrow[d] \& P_{1}\arrow[r,phantom,"\subset"]\arrow[dr] \& C\otimes C'\arrow[d,shift right=.7ex,"\alpha"']\arrow[d,shift left=.7ex,"\beta"]\\
P_{2}\otimes s\arrow[r,phantom,"\subset"]\&P_{1}\unsure{^{i}|^{\prime}}\otimes s\arrow[rr,phantom,"\subset"]\&\&C\otimes C'\otimes s
\end{tikzcd}}
\[
P
\]
\[
\deleted{\unsure{\Delta_P\otimes|A\otimes}}\quad s^{\unsure{h|k}}\otimes_{\unsure{h|k}}s^{h}=s\qquad
u\otimes1-1\otimes u\quad\unsure{\ne|\#|\times}
\]
\[
P_{\infty}=\bigcap P_i\qquad C''\subset C\otimes_{k}C'
\]
\[
P_{\infty}\deleted{\otimes s}\longrightarrow C\otimes C'\otimes\unsure{e|s|C}
\qquad \deleted{\underline{M}\otimes_{k}N}\qquad
u_C\otimes\mathrm{id}_{C'}=\mathrm{id}_{C}\otimes u_C
\]
\diagram{\begin{tikzcd}[ampersand replacement=\&]
P_\infty\arrow[dr]\&\\
\&\displaystyle\bigcap_{\unsure{i}}P_i\otimes s=P_\infty\otimes s
\end{tikzcd}}
\[
\begin{aligned}
x\otimes1&-1\otimes x\\
x\otimes y&-1\otimes xy\\
V(C)&\boxtimes_{V(s)}V(C')
\end{aligned}
\]

\uline{\unsure{Copr.|Copro.} dans \unsure{leur monde}}
\[
C\longrightarrow C\boxtimes_{s}C\qquad
V(C)\times V(C^{\prime\prime})\longrightarrow V(C'')
\]
\[
\deleted{M\otimes M'\otimes s\otimes s}\qquad V(\unsure{s|-s-|S})\quad\text{\unsure{acc. qlq}}
\]
\[
M\otimes\underline{s}\qquad M'\otimes s\qquad
\Longleftrightarrow C''\longrightarrow C\boxtimes C'\quad\text{hom de \unsure{Coalgèbres|coalgèbres}}
\]
\scanpage{284}
\mspage{- 32 -}\hfill IX
% archive sheet number: 284
% mounting strip and printed label: 350

\uline{Proposition} 3.5. (\uline{critère galoisien de réduction semi-stable}). \uline{On suppose} $S$ \uline{un trait}. \uline{Les conditions suivantes sur le schéma abélien} $A_K$ \uline{sont équivalentes} (\uline{les notations étant celles de} (2.5.3), \uline{appliquées à la situation déduite de} $(A_K,S)$ \uline{par passage au hensélisé de} $S$) :

(i) $A_K$ \uline{a réduction semi-stable sur} $S$ (3.4).

(ii) \uline{On a}, \uline{avec les notations} (2.5.3), (2.5.4) :
\[
V\handwritten{\subset}V^{\handwritten{\perp}}\quad .
\]

(iii) \uline{On a pour la partie torique} $W$ \uline{du module de Tate} :
\[
W=V^{\handwritten{\perp}}\quad .
\]

(iv) \uline{L'action du groupe d'inertie} $I$ \uline{sur} $U=T_{\handwritten{\ell}}(A(\overline{K}))$ \uline{est unipotente d'échelon} 2 (\uline{i}.\uline{e}. \uline{il existe un sous-module} $U'$ \uline{de} $U$ \uline{stable par} $I$ \uline{tel que} $I$ \uline{opère trivialement sur} $U'$ \uline{et sur} $U/U'$).

On peut supposer $S$ hensélien (3.4.0). Avec les notations de 2.1.9, (i) signifie que l'on a
\[
\handwritten{\lambda}=0\quad ,
\]
ce qui en vertu de (2.5.5) équivaut à $V=V\handwritten{\cap}V^{\handwritten{\perp}}$ i.e. à (ii). Comme $W=V\handwritten{\cap}V^{\handwritten{\perp}}$ par le théorème d'orthogonalité 2.4, on a (ii) $\handwritten{\Longleftrightarrow}$ (iii). D'ailleurs (iv) signifie aussi que $I$ opère trivialement sur $U^{I}=V$, ce qui par dualité signifie qu'il opère trivialement sur $V^{\handwritten{\perp}}$ i.e. que $V^{\handwritten{\perp}}\handwritten{\subset} U^{I}=V$, ce qui est (ii). Cela achève la démonstration.

\uline{Remarques} 3.5.1. Les conditions (ii) et (iv) ne dépendent visiblement que de l'action de $I$ sur le $\mathbb{Q}_{\handwritten{\ell}}$-vectoriel $U\otimes_{\mathbb{Z}_{\handwritten{\ell}}}\mathbb{Q}_{\handwritten{\ell}}$, et auraient pu être formulées en termes de cette action. Signalons aussi que ces conditions sont manifestement stables par tout changement de base par un morphisme
\scanpage{285}
% archive sheet number: 285
\diagram{\begin{tikzcd}[ampersand replacement=\&,column sep=large]
\unsure{s|S}\arrow[r,"\Delta_{\unsure{B|s|S}}"]\&\unsure{B|s|S}\otimes B\&\\
C\arrow[u,"\unsure{g|u}"']\arrow[r,"\Delta_C"]\&C\otimes C\arrow[u]\arrow[r,"\mathrm{id}_C\otimes g"]\&C\otimes s
\end{tikzcd}}

$V(s)$ Anneau opère sur $V(C)$

$V(C)$ structure d'algèbre sur $V(s)$ ?

\[
V(C)\mathbin{\deleted{\otimes_{\unsure{V(s)}}}}\times V(C')\xrightarrow[\text{bil. p.r. à }V(s)]{}V(C'')
\]

$1^\circ)$ bil. p.r. à $\deleted{\unsure{V(k)|V(s)}}\unsure{\mathcal O|0|\odot}$ \unsure{donc|doit|d'un} \unsure{accouplement|exprimer} :
\[
C''\longrightarrow C\otimes_{k}C'
\]
\diagram{\begin{tikzcd}[ampersand replacement=\&,column sep=huge,row sep=large]
V(C)\times V(\deleted{\unsure{C'}}s)\times V(\deleted{\unsure{s}}C')\arrow[r,"\mathrm{id}_{V(C)}\times\unsure{w|\omega}_{V(C')}"]\arrow[d,"\deleted{w_{V(C)}\times\mathrm{id}_{V(\unsure{s''|s|S})}}"']\&V(C)\times V(C')\arrow[d]\\
V(C'')\times V(\deleted{\unsure{C'}}s)\arrow[r]\&V(C'')
\end{tikzcd}}

\[
\deleted{\uparrow}\quad C''
\]
\[
\begin{aligned}
V(C)\times\operatorname{Spec}(A)&\longrightarrow V(C'')\\
V(C)&\longrightarrow V(C'')\\
\deleted{V(C)\times V(s)\times\operatorname{Spec}(A)}&
\end{aligned}
\]
\diagram{\begin{tikzcd}[ampersand replacement=\&]
C\arrow[d]\&C''\arrow[l]\arrow[d]\\
C\otimes s\&C''\otimes s\arrow[l]
\end{tikzcd}}
\[
C\otimes A\longleftarrow C''
\]
\[
\left.\begin{aligned}
V(C)\times V(C')&\mathrel{\unsure{-|\to}}V(C'')\\
\operatorname{Sym}(C)\otimes\operatorname{Sym}(C')&\longleftarrow C''
\end{aligned}\right\}
\]
\diagram{\begin{tikzcd}[ampersand replacement=\&]
C''\arrow[r,dash]\arrow[d]\&C\otimes A\arrow[d]\\
C''\otimes s\arrow[r]\&C\otimes s\otimes A
\end{tikzcd}}
\[
\deleted{C\otimes\operatorname{Sym}(C')\ \unsure{\longrightarrow y|\cap}}\qquad
\]
\[
\begin{aligned}
\Delta_C\otimes\mathrm{id}_{C'}&\simeq\mathrm{id}_{C}\otimes\Delta_{C'}\quad\unsure{\in|\subset|\swarrow}\\
&\simeq(\alpha\otimes\mathrm{id}_{s})\circ\Delta_{C''}
\end{aligned}
\]
\diagram{\begin{tikzcd}[ampersand replacement=\&,column sep=large]
C''\arrow[r,"\alpha",dash]\arrow[d,"\Delta_{C''}"']\&C\otimes C'\arrow[d,"\Delta_C\otimes\mathrm{id}_{C'}"]\arrow[dr,bend left=65]\&\\
C''\otimes s\arrow[r,"\alpha\otimes\mathrm{id}_{s}"']\&C\otimes s\otimes C'\&{}
\end{tikzcd}}
\scanpage{286}
\mspage{- 103 -}\hfill IX
% archive sheet number: 286
% mounting strip and printed label: 421

canonique de $(\overset{\handwritten{\scriptscriptstyle\vee}}{M}'\otimes_{\mathbb{Z}}\overset{\handwritten{\scriptscriptstyle\vee}}{M})\otimes_{\mathbb{Z}}\mathbb{Z}_{\handwritten{\ell}}$, d'ailleurs invariant par $\handwritten{\pi}_{\unsure{o|0}}=\handwritten{\pi}/I=\operatorname{Gal}(\overline{k}/k)$, puisque par définition il est invariant par transport de structure :
\[
u\handwritten{\in}(\overset{\handwritten{\scriptscriptstyle\vee}}{M}'\otimes_{\mathbb{Z}}\overset{\handwritten{\scriptscriptstyle\vee}}{M})\otimes_{\mathbb{Z}}\mathbb{Z}_{\handwritten{\ell}})^{I}\simeq\operatorname{Hom}_{\mathbb{Z}}(\underline{M}\otimes\underline{M}\unsure{',|;}\mathbb{Z}_{\handwritten{\ell}})\simeq\operatorname{Hom}_{\mathbb{Z}_{\handwritten{\ell}}}(\underline{M}_{\handwritten{\ell}}\otimes\underline{M}'_{\handwritten{\ell}},\mathbb{Z}_{\handwritten{\ell}})\quad .\tag*{(9.2.9)}
\]
C'est (à la notation $u$ au lieu de $u_{\handwritten{\ell}}$ près) l'élément (9.1.2) que nous voulions définir.

Pour définir (9.1.2) dans le cas général, nous aurons besoin de quelques préliminaires, que nous allons développer dans les deux sections suivantes .

9.3. \uline{Extensions panachées}.

Soient $\underline{C}$ une catégorie abélienne, et
\[
P,\ R,\ Q\tag*{(9.3.1)}
\]
trois objets de $\underline C$. On s'intéresse à la classification des objets $X$ de $\underline C$, munie d'une filtration en trois crans
\[
X^{\unsure{o|0}}=X\handwritten{\supset}X^{1}\handwritten{\supset}X^{2}\handwritten{\supset}X^{3}=0\quad ,\tag*{(*)}
\]
avec des isomorphismes donnés
\[
X^{o}/X^{1}\xrightarrow{\sim}P\quad ,\quad X^{1}/X^{2}\xrightarrow{\sim}R,\quad X^{2}/X^{3}=X^{2}\xrightarrow{\sim}Q\quad .
\]
Posant
\[
E=X^{o}/X^{2}\quad ,\quad F=X^{1}/X^{3}=X^{1}\quad ,
\]
on voit que $X$ donne naissance à deux extensions ordinaires
\[
0\longrightarrow P\longrightarrow F\longrightarrow R\longrightarrow0\quad ,\quad 0\longrightarrow R\longrightarrow E\longrightarrow Q\longrightarrow0\quad ,\tag*{(9.3.2)}
\]
de $R$ par $P$ et de $Q$ par $R$ respectivement. Notre point de vue sera ici
\scanpage{287}
% archive sheet number: 288
(1) $S$ anneau \inserted{$\mathbb Z$-}gradué \inserted{\unsure{comm|comm.}}, \uline{à degrés $\geq0$}, $\displaystyle S=\sum_{n\geq0} S_{\inserted{n}}^{\deleted{\unsure{n}}}$

\deleted{$\unsure{S_0=k}$}

(2) $k=S_{\inserted{0}}^{\deleted{\unsure{0}}}$

(3) $\operatorname{Modgrad}(S)=\mathcal M$

(4) $\underline t$ foncteur translation
\[
\left\{\begin{aligned}
&\underline t:\mathcal M\mathrel{\unsure{-|\to}}\mathcal M\\
(4.1)\quad &\underline t(M)_i=M_{i+1}
\end{aligned}\right.
\qquad\text{c'est un \uline{automorphisme} de $\mathcal M$}
\]

(4.2) \deleted{\unsure{M|$\mathcal M$}} \unsure{on pose} \inserted{$M[i]=\underline t^{\,i}(M)$.}

Pour tout entier $N$, on considère les deux sous-catégories \inserted{\unsure{str.|strict.}} pleines \inserted{$\mathcal M_N,\mathcal M^{N}$} de $\mathcal M$ \unsure{définies} par
\[
\left\{\begin{aligned}
\operatorname{Ob}\mathcal M_N&=\{P\in\operatorname{Ob}\mathcal M\mid P_i=0\quad\forall i<N\}\\
\operatorname{Ob}\mathcal M^{N}&=\{P\in\operatorname{ob}\mathcal M\mid P_i=0\quad\forall i\unsure{>|{\geq}}N\}
\end{aligned}\right.\tag*{(5)}
\]

\uline{NB} Les $\mathcal M_N$ ($\mathcal M^{N}$) décroissent (croissent) pour $N$ croissant. On a
\[
\left\{\begin{aligned}
t^{i}(\mathcal M_N)&=\mathcal M_{N-i}\\
t^{i}(\mathcal M^{N})&=\mathcal M_{\deleted{\unsure{N-i}}}^{N-i}
\end{aligned}\right.\tag*{(6)}
\]
\deleted{\unsure{donc}}

En particulier \deleted{\unsure{(6)|est}}
\[
\left\{\begin{aligned}
\deleted{\unsure{t}}\,\mathcal M_N&=t^{-N}\mathcal M_0\quad,\\
\mathcal M^{N}&=t^{-N}\mathcal M^{0}
\end{aligned}\right.\tag*{(6.1)}
\]

\inserted{Dans} la catégorie $\mathcal M$, munie déjà de l'\unsure{autom.} $t$, la donnée des familles $(\mathcal M_N)$ et $(\mathcal M^N)$ avec leurs inclusions définitoires
\[
\mathcal M_N\supset\mathcal M_{N'}\quad\text{et}\quad \mathcal M^N\subset\mathcal M^{N'}\quad\text{si}\quad N\leq N'\tag*{(7)}
\]
équivaut à celle \deleted{\unsure{des}} du couple
\[
\mathcal M_0,\mathcal M^0\quad\text{avec}\quad t\mathcal M_0\unsure{\subset|\supset}\mathcal M_0,\quad t\mathcal M_{\deleted{\unsure{0}}}^{0}\supset\mathcal M_{\deleted{\unsure{0}}}^{0}\tag*{(8.)}
\]
On va voir d'ailleurs comment $\mathcal M_{\unsure{N}},\mathcal M^{N-1}$ \deleted{(\unsure{ou encore} $\unsure{\mathcal M_N,\mathcal M^N}$)} se déterminent mutuellement -- donc $\mathcal M_0$ \uline{ou} $\mathcal M^0$ permet de \unsure{retrouver|retirer} tous \deleted{\unsure{les}} $\mathcal M_N,\mathcal M^N$.
\scanpage{288}
\mspage{- 102 -}\hfill IX
% archive sheet number: 289
% mounting strip and printed label: 420

Or par le théorème d'orthogonalité appliqué a la fois pour $A_K$ et $A'_K$ on a
\[
V^{\handwritten{\perp}}=W\quad ,\quad (U/V)^{\handwritten{\vee}}\simeq W'(-1)\quad ,
\]
où $W'$ est la partie torique de $U'=T_{\handwritten{\ell}}(A'(\overline K))$, et où le $(-1)$ désigne le twist de Tate. D'autre part, l'homomorphisme (9.2.4) se factorise par $I(\handwritten{\ell})\simeq\mathbb Z_{\handwritten{\ell}}(1)$, de sorte que l'homomorphisme (9.2.4) peut s'interprèter comme un homomorphisme canonique
\[
u:\mathbb Z_{\handwritten{\ell}}(1)\longrightarrow W'\otimes W(-1)\quad ,\tag*{(9.2.5)}
\]
ou encore comme un élément
\[
u\handwritten{\in}W'\otimes W\,(-2)\quad .\tag*{(9.2.6)}
\]

Se rappelant de la définition de $W$ et de $W'$, comme les groupes de points à valeurs dans $\overline K$ des deux faisceaux $\handwritten{\ell}$-adiques constants tordus $T_{\handwritten{\ell}}(A^{\unsure{o|0}})^{t}$ et $T_{\handwritten{\ell}}(A'^{o})^{t}$, dont les fibres spéciales sont respectivement $T_{\handwritten{\ell}}(T_o)$ et $T_{\handwritten{\ell}}(T'_o)$, on trouve les isomorphismes canoniques
\[
\left\{\begin{aligned}
W&\simeq T_{\handwritten{\ell}}(T_o)(\overline k))\simeq\overset{\handwritten{\scriptscriptstyle\vee}}{M}\otimes\mathbb Z_{\handwritten{\ell}}(1)\\
W'&\simeq T_{\handwritten{\ell}}(T_o(\overline k))\simeq M'\otimes\mathbb Z_{\handwritten{\ell}}(1)\quad ,
\end{aligned}\right.\tag*{(9.2.7)}
\]
$M$ et $M'$ désignant les groupes des caractères relativement à $\overline k$ de $T_o$ et de $T'_o$ respectivement :
\[
T_o=D(M)\quad ,\quad T'_o=D(M')\quad .\tag*{(9.2.8)}
\]
Par suite, l'élément $u$ de (9.2.6), qui exprime l'opération du groupe d'inertie $I$ sur le module de Tate $U$ de $A_K$, s'interprète comme un élément
\scanpage{289}
% archive sheet number: 290
On considère les foncteurs d'inclusion
\[
i_N:\mathcal M_N\hookrightarrow\mathcal M\qquad i_{\deleted{\unsure{N}}}^{N}:\mathcal M^{N}\mathrel{\unsure{-|\to}}\mathcal M\tag*{(9)}
\]
plus généralement, pour $N\leq N'$
\[
i_{N,N'}:\mathcal M_{N'}\hookrightarrow\mathcal M_N\qquad i^{N',N}:\mathcal M^{N}\mathrel{\unsure{\subset|\hookrightarrow}}\mathcal M^{N'}\tag*{(9.1)}
\]

\deleted{On va étudier les}

donc on a
\[
i_N\,i_{N,N'}=i_{N'\deleted{\unsure{,N}}}\quad,\quad i^{N'}i^{N',N}=i^{N}\tag*{(9.2)}
\]
\deleted{plus \unsure{g}} et si $N\leq N'\leq N''$, d'où \deleted{$\unsure{i_N}$}
\[
\left\{\begin{aligned}
\mathcal M_{N''}\subset\mathcal M_{N'}\subset\mathcal M_N\ \deleted{\text{\unsure{et}}}&\qquad\mathcal M^{N}\subset\mathcal M^{N'}\subset\mathcal M^{N''}\\
i_{N,N''}=i_{N,N'}i_{N',N''}&\qquad i^{N'',N}=i^{N'',N'}i^{\unsure{N'N|N',N}}
\end{aligned}\right.\tag*{(9.3)}
\]

On va étudier la suite de foncteurs adjoints associés à $i_N,i_{N,N'}$ d'une part, $i^N,i^{N,N'}$ d'autre part.

\uline{Proposition} 1 \quad $i_N$\inserted{$:\mathcal M_N\hookrightarrow\mathcal M$} s'insère dans une suite de 4 foncteurs adjoints
\[
q_N\unsure{,}\quad i_N\quad \tau_N\quad \varphi_N\tag*{(10)}
\]
(chaque foncteur \deleted{\unsure{étant|adjoint}} parmi $i_N,\tau_N,\varphi_N$ étant adjoint à droite de celui écrit à \unsure{sa|2} gauche) définis ainsi :

a) $\tau_N:\mathcal M\longrightarrow\mathcal M_N$ (foncteur \unsure{tronquage|troncature} à droite)
\[
\tau_{\unsure{N-1|N}}(P)\mathrel{\underset{\substack{\downarrow\\\text{hom d'adjonction}}}{\subset}}\deleted{\unsure{i}}P
\quad\text{(sous objet de $P$)}\quad.
\]
\[
\begin{aligned}
&\overset{\Vert}{i_N\tau_N(P)}\\
(\tau_N(P))_{i}&=\left\{\begin{aligned}
P_i&\quad\unsure{=}\ \text{si }i\geq N\\
0&\quad\text{si }i<N
\end{aligned}\right.
\end{aligned}\tag*{(10.1)}
\]
\scanpage{290}
\mspage{- 101 -}\hfill IX
% archive sheet number: 291
% mounting strip and printed label: 419

envisagé dans (2.5.3), et plus précisément, nous verrons que sa connaissance équivaut essentiellement à celle de l'action du groupe de monodromie $I$ sur $U$.

9.2. \uline{Cas} $\handwritten{\ell\ne}p$ . Plaçons nous donc d'abord dans le cas, plus simple, où $\handwritten{\ell\ne}p$, en reprenant alors les notations introduites dans (2.5.3). Notons qu'à torsion près, la relation $V=U^{I}$ signifie aussi que $V^{\handwritten{\perp}}$ est le module engendré par les éléments $gx-x$ ($u\handwritten{\in}U$, $g\handwritten{\in}I$), et qu'on a en tous cas
\[
gx-x\handwritten{\in}V^{\handwritten{\perp}}\quad\text{pour }x\handwritten{\in}U,\ g\handwritten{\in}I\quad .\tag*{(9.2.1)}
\]
On a donc, pour $g\handwritten{\in}I$ :
\[
gx=x+u(g).x\quad ,\tag*{(9.2.2)}
\]
où $g\handwritten{\longmapsto}u(g)$ est un homomorphisme de $U$ dans $V^{\handwritten{\perp}}$ qui s'annule sur $V=U^I$, ou ce qui revient au même
\[
u\handwritten{\in}\operatorname{Hom}(U/V\ ,V^{\handwritten{\perp}})\simeq (U/V)^{\handwritten{\vee}}\otimes V^{\handwritten{\perp}}\quad .\tag*{(9.2.3)}
\]
Dans le cas envisagé de la réduction semi-stable, i.e. quand on a $V^{\handwritten{\perp}}\handwritten{\subset} V$, la formule (9.2.2) exprime que l'opération $g$ dans $U$ est une \uline{transvection} parallèle à $V^{\handwritten{\perp}}=W$, et la composition de ces transvections n'est autre que l'addition des homomorphismes $u$ correspondants (9.2.3). Donc l'opération du groupe d'inertie $I$ sur $U$ s'explicite entièrement alors par un homomorphisme de groupes
\[
I\longrightarrow (U/V)^{\handwritten{\vee}}\otimes V^{\handwritten{\perp}}\quad .\tag*{(9.2.4)}
\]
\margin{\handwritten{$\unsure{\longleftarrow}$}}
\scanpage{291}
% archive sheet number: 292
b) $\varphi_N:\mathcal M_N\longrightarrow\mathcal M$
\[
\left|\begin{aligned}
(\varphi_N(Q))_i&=\operatorname{Hom}_{\mathcal M}^{i}(\tau_{N-i}(S),\overset{\substack{\inserted{i_N Q}\\\Vert}}{Q})(\longrightarrow\operatorname{Hom}_{\mathcal M}^{0}(S,Q)=Q_{\unsure{i}}\unsure{)}\\
&=\operatorname{Hom}_{\mathcal M}^{\deleted{\unsure{i|*}}}(\tau_{N-i}(S)[-i],Q)
\end{aligned}\right.\tag*{(10.2)}
\]
Avec \deleted{\unsure{leur}} structure de module provenant (pour $k\in\mathbb Z,\ u\in S_{\inserted{k}}^{\deleted{\unsure{k|*}}}$) des opérations, multiplication par $u$ :
\[
\tau_{N-i}(S)\xleftarrow{\deg k}\tau_{N-i-k}(S)
\]
en transposant ; on a un hom \deleted{dans} \deleted{\unsure{mod}} $\mathcal M$ décrit dans (10.2)
\[
\varphi_N(Q)\longrightarrow i_N(Q)\tag*{(10.2.1)}
\]
d'où en appliquant $\tau_N$
\[
\tau_N\varphi_N Q\xrightarrow[\text{adj.}]{\sim}Q\tag*{(10.2.2)}
\]
qui est un \unsure{isom|isom.} (c'est un hom d'adjonction)

L'autre hom d'adjonction
\[
\deleted{\varphi_N\tau_N P\longleftarrow}P\longrightarrow\varphi_N\tau_N P
\]
est déduit des \deleted{\unsure{h'}\ \unsure{d'adjonction}} \inserted{hom}
\[
\deleted{\tau_N P\longrightarrow P}\quad \operatorname{Hom}_{\unsure{\mathcal M}}^{i}(\tau_{N\deleted{\unsure{-i}}}S,P)
\]
\noindent (10.2.3)
\diagram{\begin{tikzcd}[ampersand replacement=\&,column sep=small,row sep=large]
P_i\arrow[rr]\arrow[dr,"\unsure{\sim|2|\simeq}"']\&\&(\varphi_N\tau_N P)_i\simeq\operatorname{Hom}^{i}(\tau_{N-i}S,\tau_N P)\\
\&\operatorname{Hom}^{i}(S,P)\arrow[ur]\&
\end{tikzcd}}
\margin{tout hom de degré $i$ \deleted{de $S$ dans $P$} $\tau_{N-i}S$ dans $P$ est à valeurs dans $\tau_{N\deleted{\unsure{-i}}}P\subset P$}

c) (NB On n'utilisera pas $q_N$, on le met par acquit de conscience) :
\[
q_N(P)=P/\text{sous-module engendré par les }P_i\text{ pour }i<N
\]
\scanpage{292}
% archive sheet number: 293
et que ces suites et la suite constante de valeur 1 forment une base de l\deleted{\unsure{a}}'espace vectoriel $E_r$ de toutes les suites complexes de période $r$. Montrer que la série
\[
u(s)_n c_n
\]
est alternée à valeurs absolues décroissantes tendant vers zéro, donc convergente. En déduire une autre démonstration de b).

d) Soit $z$ un nombre complexe, et posons $u_n=\unsure{z|Z}^{n}$ pour $n\quad 1$. \deleted{\unsure{D}}Montrer que la suite des sommes partielles $S_n$ de (1) reste bornée si et seulement si on a
\[
z\quad1\ ,\quad z\ne1\tag*{(5)}
\]
(auquel cas la suite des $S_n$ est majorée par $\;2\,\dfrac{z}{1-z}\;$). \deleted{Montrer que si la série $c_n$ est divergente, alors} \inserted{En déduire que} la condition (5) implique que la série \deleted{$\unsure{u_n c_n|c_n}$} \inserted{$c_n z^n$} est convergente. Montrer qu'elle est même absolument convergente pour $z\quad1$, tandis que pour un $z$ tel que $\;z\;=1$, $z\ne1$, on a convergence absolue si et seulement si la \deleted{suite} série $c_n$ est convergente. Enfin, \deleted{\unsure{Application} à la convergence et à la convergence absolue\unsure{:} montrer que pour toute partie $A$} \inserted{$A$ désignant une partie} de l'ensemble \inserted{$E$} des nombres complexes satisfaisant (5), pour que la convergence de la série $c_n z^n$ soit uniforme pour $z\quad A$, il suffit qu'il existe un nombre $a\quad0$ tel que $z\quad A$ implique $1-z\quad a$. Montrer que \deleted{\unsure{c}}cette condition est aussi nécessaire si la série $c_n$ est divergente (tandis que si $c_n$ est convergente, la convergence est uniforme pour l'ensemble des nombres complexes $z$ tels que $\;z\quad1$). \deleted{\unsure{E}}

e) Montrer que la fonction $f(z)$ sur l'ensemble $E$ des nombres complexes $z$ satisfaisant (5), représentée par la série
\[
f(z)=c_n z^n\tag*{(6)}
\]
est continue. Elle admet une limite finie pour $z=1$ si et seulement si la \deleted{suite} série $c_n$ est convergente, auquel cas la fonction $f(z)$ définie par la formule (6) est définie et continue \deleted{pour} sur tout le disque unité $z\quad1$.

f) Appliquer les résultats de d) et e) à la série trigonométrique
\[
c_n e^{ix}\quad ,
\]
et aux séries
\[
c_n\cos nx\ ,\quad c_n\sin nx\quad :
\]
en particulier, ces séries sont convergentes pour tout $x$ réel n'appartenant pas à l'ensemble \deleted{$\unsure{2\ \mathrm R|2\ \mathrm Z}$} $2\quad\underline{\mathrm Z}$ (i.e. $x$ pas de la forme $2\ n$, avec $n$ entier), et représentent des fonctions continues de $x$ pour $x\quad\underline{\mathrm R}-2\quad\underline{\mathrm Z}$ .

g) Montrer que la suite $(z^n)_{n\quad1}$ ($z$ étant un nombre complexe) est périodique de période $r$ si et seulement si $z$ est une racine r.ème de l'unité. \deleted{\unsure{Lorsque}} Déduire de ceci et et b) que lorsque $z$ est une racine de l'unité la série \deleted{\unsure{En}} $c\ z^n$ est convergente, \unsure{dauf|sauf} dans le cas où $z=1$ et $c_n$ \unsure{est}
\scanpage{293}
% archive sheet number: 294
% Vertical stroke to the left of the corollary.
\uline{Cor} 1

\noindent\begin{tabular}{|p{.91\linewidth}}
Les foncteurs $i_N,\tau_N$ commutent à \unsure{ttes|toutes} les sortes de limites inductives et projectives, $q_N$ commute aux $\varinjlim$, $\varphi_N$ aux $\varprojlim$.
Les foncteurs $i_N,\varphi_N$ sont pl. fid, les foncteurs $q_N,\deleted{\unsure{\varphi}}\tau_N$ sont des foncteurs localisation. \deleted{\unsure{Comme}} Pour $\tau_N$ (\unsure{puisque|comme} il est exact, et commute aux $\varinjlim$ et $\varprojlim$) c'est la localisation mod une sous-cat. abélienne épaisse, stable par $\varinjlim$ et $\varprojlim$ ; en fait, c'est précisément la sous-cat $\mathcal M^{N-1}$ (qui est ainsi décrite \deleted{\unsure{via}} en termes de $\mathcal M_N$).
\end{tabular}
\scanpage{294}
% archive sheet number: 295
Problème MP 2

Un des objets du présent exercice est d'étudier la convergence de certaines séries de la forme
\[
u_n/\unsure{n|m}\quad ,
\]
où $(u_n)$ est une suite bornée de nombres complexes, et d'en calculer la somme dans certains cas, notamment dans le cas où la suite $(u_n)$ admet une ``période'' $r$ (de sorte qu'on a $u_{n+r}=u_n$ pour tout $r$).

I

Soit $(c_n)_{n\quad1}$ une suite décroissante de nombres positifs tendant vers \deleted{0, et soit \unsure{(u}} 0, et soit $(u_n)_{n\quad1}$ une suite de nombres complexes. \inserted{a)} On suppose qu'il existe un nombre $M\quad0$ tel que pour tout entier $n\quad0$, on a
\[
S_n={}_{1\ i\ n}\;u_i\quad \unsure{.}\quad M\quad .
\]
\deleted{a)} Prouver que la série
\[
u_n c_n
\]
est convergente, et que pour tout entier $n\quad1$, le reste
\[
R_n={}_{i\ \,n+1}\;u_i c_i
\]
satisfait
\[
R_n\quad Mc_{n+1}\quad .
\]
(Utiliser le critère de Cauchy, en prouvant que pour $m\quad n$, on a
\[
{}_{n\ i\ m}\;u_i c_i\qquad Mc_{n+1}\quad ;
\]
pour ceci, remplacer dans le premier membre $c_i$ par $S_i-S_{i-1}$.)

b) On suppose que la suite $(u_n)$ est périodique de période $r$ ($r\quad1$), et on s'intéresse à la convergence de la série $\;u_n c_n$. Montrer que si la série $\;c_n$ est convergente, il en est de même de $\;u_n c_n$ \deleted{\unsure{q|u}}quelle que soit la suite (périodique) $(u_n)$. Montrer que si la série $\;c_n$ est divergente, alors $\;u_n c_n$ est convergente si et seulement si on a
\[
{}_{1\ i\ r}\;u_i=0\unsure{/}\quad .
\]
(Montrer que toute suite $(u_n)$ périodique de période $r$ se met de façon unique sous forme d'une somme de deux suites périodiques de période $r$
\[
u_n=u'_n+u''_n\quad ,
\]
où $(u'_n)$ est une suite constante et où $(u''_n)$ satisfait à la condition
\[
{}_{1\ i\ r}\;u''_i=0\quad .
\]
Appliquer alors le résultat de a) à cette dernière suite.)

c) On considère l'espace \inserted{vectoriel\unsure{:}complexe $E_r$} des suites périodiques \inserted{$u=$}$(u_n)$ de période $r$ satisfaisant la condition (3). Soit $n_{\unsure{o|0}}\quad1$, et soit , pour tout entier $\deleted{\unsure{i}}s$ tel que $1\quad\deleted{\unsure{i}}s\quad r-1$, $u(\deleted{\unsure{i}}s)=(u(\deleted{\unsure{i}}s)_n)_{n\quad1}$ l'unique suite $\;E_r$ telle qu'on ait
\[
\begin{aligned}
\text{pour }0\quad i\quad r-1,\quad \deleted{u(i)_{\unsure{n_o+i}}}\;u(\deleted{\unsure{i}}s)_{n_o+\unsure{i|1}}&=\begin{aligned}
1&\ \text{si }\deleted{\text{\unsure{$i$ est 0, ou $s=1$ si}}}\ i=0\\
-1&\ \text{si }i=s\\
0&\ \text{si }i\ne0,s
\end{aligned}\quad .
\end{aligned}
\]
Prouver que les $u(s)$ ($1\quad s\quad r-1$) forment une base de l'espace vectoriel
\scanpage{295}
% archive sheet number: 296
$\unsure{\mathcal C|C}$ cat. abélienne

$S$ anneau $\mathbb Z$-gradué comm. à degrés $\geq0$

$\mathcal M$\inserted{$=\mathcal M(\mathcal C,S)$} catégorie des $\mathcal C$-$S$-modules \uline{gradués}.

$\mathcal M_N,\ \mathcal M^N$ \unsure{définis} comme plus haut.
\[
\mathcal M_N\xrightarrow{i_N}\mathcal M
\]
admet foncteur adjoint à droite
\[
\begin{aligned}
\mathcal M&\xrightarrow{\tau_N}\mathcal M_N\\
i_N\tau_N(P)&\subset P
\end{aligned}
\]
\deleted{qui admet}

Construisons adjoint à droite $\varphi_N$ de $\tau_N$ :
\[
\operatorname{Hom}_{\mathcal M_N}(\tau_N P,\underset{\substack{\rotatebox{-90}{$\in$}\\\mathcal M_N}}{Q})\simeq
\operatorname{Hom}_{\mathcal M}(P,\varphi_N Q)\quad ?\tag*{\unsure{1}}
\]
Il faut donc prouver que pour $Q$ fixé, le premier \unsure{mb}, comme foncteur de $P$, est repr. Or il est clair que $\tau_N$ commute aux $\varinjlim$ qui sont \unsure{représent.ble} dans $\mathcal M$ donc le foncteur du premier \unsure{mb} transforme les $\varinjlim$ en $\varprojlim$. Cela implique donc ce qu'on veut lorsque $\mathcal C$ \uline{ou} $\mathcal C^{o}$ est une ``cat. de \unsure{Groth|Groth.}'' \ldots

Comment \inserted{alors} décrire les $(\varphi_N Q)_i\in\operatorname{Ob}\mathcal C$

On a pour \deleted{\unsure{$\varphi$}} tout $M\in\operatorname{Ob}\mathcal M$
\[
M_i\simeq\underline{\operatorname{Hom}}^{*}(S,M)_i
\]
donc
\[
\begin{aligned}
(\varphi_N Q)_i&=\underline{\operatorname{Hom}}^{*}(S,\varphi_N Q)_i\simeq\underline{\operatorname{Hom}}^{*}(S[-i],\varphi_N Q)_{\unsure{0}}\\
&\overset{\text{\inserted{\unsure{justifier}}}}{=}\underline{\operatorname{Hom}}^{*}(\tau_N(S[-i]),Q)_0\\
&=(\underline{\operatorname{Hom}}^{*}((\tau_{N-i}(S))[-i],Q)\deleted{\unsure{[-i]}})_0
\end{aligned}
\]
\scanpage{296}
\mspage{- 100 -}\hfill IX

9. \uline{Définition de l'accouplement de monodromie} $\underline M_{\handwritten{\ell}}\otimes\underline M'_{\handwritten{\ell}}\longrightarrow\mathbb Z_{\handwritten{\ell} S}$.

9.1. Dans le présent numéro on suppose $S$ hensélien, et que $A_K$ ait réduction semi-stable sur $S$, ou ce qui revient au même, que $A'_K$ ait réduction semi-stable sur $S$ (3.5.1). On dispose alors sur $S$ de deux $\mathbb Z$-Modules constants tordus $\underline M$ et $M'$ (5.5.2), qui sur un revêtement fini étale convenable de $S$ (donc sur $S$ lui-même si $S$ est strictement local i.e. si $k$ est séparablement clos) se déploient en des groupes constants de valeur $M$ resp. $M'$ isomorphes (non canoniquement) à $\mathbb Z^r$, $r$ étant la dimension commune des tores maximaux de $A_o$ et de $A'_o$. Nous allons considérer, pour tout nombre premier $\handwritten{\ell}$, les faisceaux $\handwritten{\ell}$-adiques correspondants
\[
\tag*{(9.1.1)}
\underline M_{\handwritten{\ell}}=\underline M\otimes_{\mathbb Z}\mathbb Z_{\handwritten{\ell}},\quad
\underline M'_{\handwritten{\ell}}=\underline M'\otimes_{\mathbb Z}\mathbb Z_{\handwritten{\ell}},
\]
dont la signification géométrique est explicitée par les formules (5.5.7) et (5.5.8). Dans le présent numéro, nous allons définir un accouplement canonique (dit ``\uline{accouplement de monodromie}'')
\[
\tag*{(9.1.2)}
u_{\handwritten{\ell}}:\underline M_{\handwritten{\ell}}\otimes\underline M'_{\handwritten{\ell}}\longrightarrow\mathbb Z_{\handwritten{\ell}},
\]
qui sera défini en termes des pro-$\handwritten{\ell}$-groupes de Barsotti-Tate $T_{\handwritten{\ell}}(A^o)$ et $T_{\handwritten{\ell}}(A'^{o})$, et même de la fibre générique $T_{\handwritten{\ell}}(A_K)$ du premier de ces groupes (sous réserve, lorsque $\handwritten{\ell}=p$, de disposer du théorème de Tate [33, th.4], i.e. pour l'instant qu'on soit alors dans le cas d'inégales caractéristiques). Lorsque $\handwritten{\ell\ne}p$, il revient au même de dire que l'accouplement (9.1.2) est décrit en termes du $\handwritten{\pi}$-module galoisien
\[
\tag*{(9.1.3)}
U=T_{\handwritten{\ell}}(A_K)(\bar K)=T_{\handwritten{\ell}}(A_{\bar K})
\]
% Two isolated dark traces below the last formula, possibly transferred.
\margin{\handwritten{$\unsure{\ni}$\qquad $\unsure{u}$}}
% mount label: 418
% archive sheet number: 297
